% Shared preamble of main.tex and supplement.tex. Load after the document class.
\usepackage{amsmath,amssymb,amsthm,mathtools}
\allowdisplaybreaks[1]
\usepackage[margin=1.1in]{geometry}
\setlength{\emergencystretch}{1em}
\usepackage{booktabs,longtable}
\usepackage{xr-hyper}
\usepackage[hidelinks]{hyperref}
\usepackage{refcount}
\usepackage{enumitem}

% ---------------------------------------------------------------- theorem environments
\newtheorem{theorem}{Theorem}[section]
\newtheorem{proposition}[theorem]{Proposition}
\newtheorem{lemma}[theorem]{Lemma}
\newtheorem{corollary}[theorem]{Corollary}
% Lettered results A--G of the introduction.
\newtheorem{maintheorem}{Theorem}
\newtheorem{mainproposition}[maintheorem]{Proposition}
\newtheorem{maincorollary}[maintheorem]{Corollary}
\renewcommand{\themaintheorem}{\Alph{maintheorem}}
\theoremstyle{definition}
\newtheorem{definition}[theorem]{Definition}
\newtheorem{construction}[theorem]{Construction}
\theoremstyle{remark}
\newtheorem{remark}[theorem]{Remark}

% ---------------------------------------------------------------- notation
\newcommand{\R}{\mathbb{R}}
\newcommand{\C}{\mathbb{C}}
\newcommand{\Z}{\mathbb{Z}}
\newcommand{\Q}{\mathbb{Q}}
\newcommand{\N}{\mathbb{N}}
\newcommand{\CP}{\mathbb{CP}}
\newcommand{\eps}{\varepsilon}
\newcommand{\Rcal}{\mathcal{R}}
\newcommand{\Hb}{H_{\mathrm b}}            % base area profile
\newcommand{\Pb}{P_{\mathrm b}}            % unused area
\newcommand{\Bo}{B^{\circ}}                % open capacity ball
\newcommand{\EB}{\mathbb{B}}               % Euclidean ball (radius)
\newcommand{\Pack}{\mathrm{Pack}}          % packing specification
\newcommand{\Hfin}{H_{\mathrm{fin}}}
\newcommand{\defdesc}{\operatorname{the}}   % definite description (never \the, a TeX primitive)
\DeclareMathOperator{\vol}{vol}

% ---------------------------------------------------------------- Lean / blueprint annotations
% Unqualified \lean{...} records references without a claim of complete coverage.
% Optional [full] / [needed] classify coverage in the PDF and generated paper index.
% These are editorial scope annotations, not a machine check of paper/Lean equivalence.
\newcommand{\leancoveragefull}{full statement verified}
\newcommand{\leancoverageneeded}{needed consequence verified}
\newcommand{\lean}[2][]{\if\relax\detokenize{#1}\relax\else
  \ifcsname leancoverage#1\endcsname\else
    \PackageError{paper}{Unknown Lean coverage `#1'}{Use full or needed.}\fi
  \par\nobreak\noindent{\footnotesize\upshape[Lean: \csname leancoverage#1\endcsname.]\par}\nobreak
  \fi\ignorespaces}
% Main-result counterparts agree up to the presentation differences in Section~\ref{sec:ver}.
\newcommand{\leanmark}[1]{\par\nobreak\noindent
  {\footnotesize\upshape\raggedright[Lean: full statement verified; \texttt{#1}]\par}\nobreak}
\newcommand{\uses}[1]{}
\newcommand{\proves}[1]{}

% ---------------------------------------------------------------- cross-document references
% Numbers of the other document carry a prefix: S for the supplement, M for the main text.
% In the main text: \Sref{Lemma}{lem:shd-foo} -> "Lemma S6.3"; \Ssec{sec:shd} -> "\S S6".
% In the supplement: \Mref{Lemma}{lem:eff-gronwall} -> "Lemma M7.1"; the lettered results of
% the main text keep their letters: \Mref{Theorem}{thm:main} -> "Theorem C".
% The second argument is the label WITHOUT the S-/M- prefix.
\makeatletter
\DeclareRobustCommand{\Mprefix}[1]{\edef\@Mref{\getrefnumber{M-#1}}\expandafter\@Mprefix\@Mref\@nil}
\def\@Mprefix#1#2\@nil{\ifnum`#1<`A M\fi}
\makeatother
\newcommand{\Snum}[1]{\href[page=\getpagerefnumber{S-#1}]{supplement.pdf}{S\ref*{S-#1}}}
\newcommand{\Mnum}[1]{\href[page=\getpagerefnumber{M-#1}]{main.pdf}{\Mprefix{#1}\ref*{M-#1}}}
\newcommand{\Sref}[2]{#1~\Snum{#2}}
\newcommand{\Ssec}[1]{\S\Snum{#1}}
\newcommand{\Mref}[2]{#1~\Mnum{#2}}
\newcommand{\Msec}[1]{\S\Mnum{#1}}
